\chapter{Embedding problems and regular realizations}
\label{ch:inverse-galois}

An inverse Galois question asks for a Galois extension with a specified
group. An embedding problem also prescribes a quotient extension and
the map to its Galois group. We first treat finite embedding problems,
including both the obstruction to a lift and the extra requirement for
surjectivity. We then pass from rational-function fields to residue
fields by specialization at an explicitly good point.

We use the finite and infinite Galois correspondences of
the earlier field-theory volume~\cite[Results F01, F07, T02,
C04, C06, C08--C10]{mg-aft}. In particular, the absolute Galois
group \(\Gamma_k=\operatorname{Gal}(k^{\mathrm{s}}/k)\) carries
the Krull topology, and finite groups have the discrete topology.

\section{Finite embedding problems}

\begin{definition}[Weak and proper solutions]
\label{res:IG01}
A finite embedding problem for a profinite group \(\Gamma\)
consists of continuous surjections
\[
   \rho\colon\Gamma\twoheadrightarrow Q,\qquad
   \pi\colon B\twoheadrightarrow Q
\]
with \(Q\) and \(B\) finite. Write \(A=\ker\pi\). A weak solution is
a continuous homomorphism \(\lambda\colon\Gamma\to B\) satisfying
\(\pi\lambda=\rho\). It is proper if it is surjective. The problem is
split if \(\pi\) has a homomorphic section and central if \(A\) is
contained in the centre of \(B\).
\end{definition}

\begin{proposition}[The prescribed quotient extension]
\label{res:IG02}
Let \(\Gamma=\Gamma_k\), and let \(\rho\) be restriction to a prescribed
finite Galois extension \(E/k\) in \(k^{\mathrm{s}}\), with the
identification \(Q=\operatorname{Gal}(E/k)\). A weak solution produces
a finite Galois extension \(L/k\) containing \(E\) and an injection
\(\iota\colon\operatorname{Gal}(L/k)\hookrightarrow B\) for which
\(\pi\iota\) is restriction to \(E\). Conversely such an extension
and injection give a weak solution. This is a proper solution exactly
when \(\iota\) is an isomorphism. In particular, when \(Q=1\),
proper solutions are finite Galois realizations of \(B\) over \(k\).
\end{proposition}

\begin{proof}
The kernel of \(\lambda\) is open and normal and is contained in
\(\ker\rho\). The fixed field \(L\) of this kernel is finite and
Galois over \(k\) and contains \(E\). Restriction gives the
isomorphism from the quotient \(\Gamma_k/\ker\lambda\) to
\(\operatorname{Gal}(L/k)\), and \(\lambda\) identifies this
quotient with its image in \(B\). The equation
\(\pi\lambda=\rho\) checks the restriction diagram. Conversely,
compose the surjective restriction to \(L\) with \(\iota\).
\end{proof}

\section{Abelian kernels and properness}

In this section cochains with values in an abelian kernel are written
additively. The coefficients in \(H^i_{\mathrm{cts}}\) have the
discrete topology. This agrees with the cochain convention in
\cite[Results GC01 and GC04--GC05]{mg-groups}.

\begin{lemma}[The class of an abelian extension]
\label{res:IG03}
For an abelian kernel \(A\) in Definition~\ref{res:IG01}, conjugation
by representatives in \(B\) defines an action of \(Q\) on \(A\).
Let \(s\colon Q\to B\) be a set section with \(s(1)=1\).
The factor set \(c(q,r)=s(q)s(r)s(qr)^{-1}\), valued in \(A\),
is a normalized \(2\)-cocycle for this action. Its class
\([c]\in H^2(Q,A)\) is independent of \(s\), and it vanishes
exactly when \(\pi\) splits.
\end{lemma}

\begin{proof}
Representatives of \(q\) differ by an element of \(A\), whose
conjugation on \(A\) is trivial. Also \(s(q)s(r)\) represents \(qr\),
so the resulting conjugations give an action. In the coordinates
\((a,q)\leftrightarrow a\,s(q)\), multiplication has the rule
\[
 (a,q)(b,r)=(a+q\cdot b+c(q,r),\,qr).
\]
Associativity implies
\[
 c(q,r)+c(qr,u)=q\cdot c(r,u)+c(q,ru).
\]
Normalization follows from \(s(1)=1\). Replacing \(s(q)\) by
\(b(q)s(q)\) adds the coboundary
\(\delta b(q,r)=b(q)+q\cdot b(r)-b(qr)\) to \(c\).
If \(c\) is a coboundary, changing the section by the negative of a
cochain representing it makes the new factor set zero. A section
with zero factor set is a homomorphism.
\end{proof}

\begin{theorem}[Obstruction and classification of weak lifts]
\label{res:IG04}
Give \(A\) the \(\Gamma\)-action through \(\rho\).
For an abelian kernel, a weak solution exists exactly when
\[
        \rho^*[c]=0\quad\text{in }H^2_{\mathrm{cts}}(\Gamma,A).
\]
If a weak solution exists, the set of weak solutions modulo
conjugation by \(A\) has a free transitive action of
\(H^1_{\mathrm{cts}}(\Gamma,A)\).
For a central kernel the action on \(A\) is trivial.
\end{theorem}

\begin{proof}
A continuous map lifting \(\rho\) has the unique form
\(\lambda(\gamma)=z(\gamma)s(\rho\gamma)\) with a continuous
map \(z\colon\Gamma\to A\). By Lemma~\ref{res:IG03}, it is a
homomorphism exactly when
\[
   z(\gamma)+\gamma\cdot z(\eta)-z(\gamma\eta)
           =-c(\rho\gamma,\rho\eta).
\]
The equation at \((1,1)\) forces \(z(1)=0\).
This gives the obstruction. Once one such \(z_0\) is fixed, all
solutions are \(z_0+w\) for continuous \(1\)-cocycles \(w\).
Conjugation by \(a\in A\) adds the cochain
\(\gamma\mapsto a-\gamma\cdot a\) to \(z\); these cochains are
exactly the \(1\)-coboundaries. This proves the claimed action.
\end{proof}

\begin{corollary}[Detection on finite stages]
\label{res:IG05}
Let \(U\) range over the open normal subgroups of \(\Gamma\) contained
in \(\ker\rho\), and write \(\rho_U\colon\Gamma/U\to Q\) for
the induced map. For abelian kernel \(A\), a weak solution exists
exactly when \(\rho_U^*[c]\) vanishes at one finite stage
\(H^2(\Gamma/U,A)\). Inflation along
\(\Gamma/V\to\Gamma/U\), for \(V\leq U\), carries the class
for \(U\) to the class for \(V\). In an absolute Galois group these
are the stages obtained by restriction to the finite Galois
extensions containing \(E\).
\end{corollary}

\begin{proof}
The proof of Theorem~\ref{res:IG04} applied to the finite group
\(\Gamma/U\) gives the equivalence between vanishing at this
stage and a weak lift from this stage. Inflation acts by
precomposition on each cochain, giving the asserted transition.
Every continuous lift to the finite group \(B\) factors through
the finite quotient by its open normal kernel, which is contained
in \(\ker\rho\). For \(\Gamma_k\), open normal kernels are exactly
the kernels of restriction to finite Galois subextensions, by
the earlier correspondence~\cite[Results C06 and C08--C09]{mg-aft}.
\end{proof}

\begin{proposition}[Proper solutions of split problems]
\label{res:IG06}
Put \(N=\ker\rho\). A weak solution of any finite embedding problem
is proper if and only if \(\lambda(N)=A\).
Suppose \(\pi\) splits and identify \(B=A\rtimes Q\) using a
homomorphic section. Weak solutions are the maps
\[
 \lambda(\gamma)=(z(\gamma),\rho\gamma),\qquad
 z(\gamma\eta)=z(\gamma)\bigl(\rho\gamma\cdot z(\eta)\bigr),
\]
for continuous \(z\colon\Gamma\to A\); \(A\) may be nonabelian
in this assertion. The restriction \(z|_N\) is a homomorphism,
and this weak solution is proper exactly when \(z(N)=A\).
\end{proposition}

\begin{proof}
The image \(H\) of any weak solution maps onto \(Q\).
The equality \(H\cap A=\lambda(N)\) follows by applying \(\pi\)
and using \(\pi\lambda=\rho\). If \(H\cap A=A\), then for
any \(b\in B\) choose \(h\in H\) over \(\pi(b)\) and write
\(b=(bh^{-1})h\in H\). The converse is immediate.
For a split problem the displayed equation for \(z\) is exactly
the multiplication condition in the semidirect product. On \(N\),
the action through \(\rho\) is trivial and \(z\) is the
restriction of \(\lambda\) with values in \(A\).
\end{proof}

\begin{example}[A central obstruction and a split weak solution]
\label{res:IG07}
Take \(\Gamma=Q=C_2\) and \(\rho\) the identity. The epimorphism
\(C_4\to C_2\) has no weak solution. Indeed the two possible
images of the generator both have order four. The lift from the
set section has \(c(\bar 1,\bar 1)=2\in\{0,2\}\subset C_4\).
For a normalized cochain \(z\) with values in this trivial-action
kernel, \(\delta z(\bar 1,\bar 1)=2z(\bar 1)=0\). Thus the
factor set represents a nonzero obstruction in the notation of
Theorem~\ref{res:IG04}. In comparison, projection
\(C_2\times C_2\to C_2\) admits the section in the first
factor. Its embedding problem has weak solutions, and
Proposition~\ref{res:IG06} rules out a proper one since \(N=1\).
\end{example}

\begin{proposition}[Finite-field embedding problems]
\label{res:IG08}
Every finite embedding problem for the absolute Galois group of
a finite field has a weak solution. It has a proper solution
exactly when \(B\) is cyclic.
\end{proposition}

\begin{proof}
The absolute Galois group is \(\widehat{\mathbb Z}\) with
arithmetic Frobenius the topological generator
\cite[Results E04, E06 and E07]{mg-aft}. The element
\(q=\rho(1)\) therefore generates \(Q\). Any \(b\in\pi^{-1}(q)\)
of order \(m\) defines a weak lift by reducing
\(\widehat{\mathbb Z}\) modulo \(m\) and sending \(1\) to \(b\).
The composite with \(\pi\) agrees with \(\rho\) on the dense
integer subgroup and hence everywhere. Since the image of
\(\widehat{\mathbb Z}\) in any finite group is generated by
the image of \(1\), properness requires \(B\) cyclic.

Conversely, suppose \(B=\langle x\rangle\) has order \(m\) and
\(Q\) has order \(n\mid m\). We have
\(q=\pi(x)^a\) for some \(a\) prime to \(n\).
Choose \(a'\equiv a\pmod n\) and \(a'\equiv1\pmod p\) for
each prime \(p\mid m\) with \(p\nmid n\), using the Chinese
remainder construction. Then \(a'\) is prime to \(m\), and
\(x^{a'}\) is a generator of \(B\) lifting \(q\).
\end{proof}

\section{Specialization at a good point}

Let \(k\) be any field and \(u\) an indeterminate. For a field
\(\ell\) containing \(k\), write \(v_u\) for the order of a nonzero
Laurent series over \(\ell\). The leading terms of nonzero
series multiply, so orders add under multiplication, and the order of
a sum is at least the lesser order. The residue of a series of
nonnegative order is its constant term.

\begin{lemma}[Simple-root lifting in power series]
\label{res:IG09}
If \(f\in k[[u]][X]\), \(\ell/k\) is a field extension, and
\(r\in\ell\) is a simple root of the reduction \(\overline f\),
there is exactly one \(\alpha\in\ell[[u]]\) with
\(\alpha(0)=r\) and \(f(\alpha)=0\).
\end{lemma}

\begin{proof}
Begin with \(\alpha_1=r\), which solves the equation modulo \(u\).
Given a polynomial \(\alpha_j\equiv r\pmod u\) with
\(f(\alpha_j)\equiv0\pmod{u^j}\) for \(j\geq1\),
every extension of its coefficients through degree \(j\) has
the form \(\alpha_j+b u^j\) modulo \(u^{j+1}\).
Expansion gives
\[
 f(\alpha_j+b u^j)\equiv f(\alpha_j)+b u^j\overline f'(r)
          \pmod{u^{j+1}},
\]
since \(2j\geq j+1\). The coefficient \(b\in\ell\) that kills the
degree-\(j\) term is unique because \(\overline f'(r)\ne0\).
The successive coefficients define a series \(\alpha\).
Its evaluation in \(f\) is zero modulo every \(u^j\), so it
is zero. The same coefficient recursion proves uniqueness.
\end{proof}

\begin{theorem}[Good specialization as a decomposition subgroup]
\label{res:IG10}
Let \(f(t,X)\in k(t)[X]\) be monic, with coefficients regular at
\(t=a\in k\), and assume \(f(a,X)\) is separable. Let \(L\) be a
splitting field over \(k(t)\) and \(\ell\) a splitting field of
\(f(a,X)\) over \(k\). There is an embedding
\(L\hookrightarrow\ell((u))\) extending \(t\mapsto a+u\)
which matches the generic roots with their distinct reductions.
For the valuation \(w\) induced by \(v_u\), its residue field is
\(\ell\), its value group is \(\mathbb Z\) with \(w(t-a)=1\),
and the map
\[
 \{\sigma\in\operatorname{Gal}(L/k(t)):
           w(\sigma x)=w(x)\ \text{for every }x\in L^\times\}
       \longrightarrow \operatorname{Gal}(\ell/k)
\]
given by reduction is an isomorphism. In particular the subgroup
stabilizing \(w\) has trivial inertia kernel, and the specialized
Galois group embeds into the generic Galois group up to conjugacy.
\end{theorem}

\begin{proof}
The substitution \(t=a+u\) gives the indicated embedding of
\(k(t)\) into \(k((u))\). A coefficient regular at \(a\) can be
written with a denominator nonzero at \(a\), whose image in
\(k[[u]]\) has an inverse by coefficient recursion.
By Lemma~\ref{res:IG09}, each root of \(f(a,X)\) lifts to a root
of \(f(a+u,X)\) in \(\ell[[u]]\). The lifted roots are distinct
by their reductions. Since the polynomial is monic of unchanged
degree, the lifts are all of its roots. They generate a splitting
field over the embedded \(k(t)\); uniqueness of the splitting field
gives the required embedding. Write \(L\) for its image.

The residue field of \(L\) lies in \(\ell\) and contains every
specialized root, so equals \(\ell\). The value group lies in
\(\mathbb Z\) and contains \(1=w(u)\). An automorphism preserving
\(w\) acts on the residue field. If it induces the identity, it
fixes the reduction of each lifted root. The reductions are distinct,
and the automorphism permutes these roots, so it fixes all of them
and is the identity. For surjectivity, extend any
\(\tau\in\operatorname{Gal}(\ell/k)\) coefficientwise to
\(\ell((u))\), fixing \(u\). This automorphism preserves order.
It takes the lift of \(r\) to the lift of \(\tau(r)\) by uniqueness,
and hence preserves \(L\). Its restriction induces \(\tau\) on
the residue field. Finally, a different embedding into
\(\ell((u))\) over the same \(t\mapsto a+u\) has the same image,
since its roots must be the same full set of lifted roots.
Thus it changes the induced valuation by an automorphism of \(L\),
which conjugates the displayed subgroup.
\end{proof}

\begin{corollary}[Specialization in a tower]
\label{res:IG11}
In Theorem~\ref{res:IG10}, let \(L_0\subseteq L\) be a splitting
field over \(k(t)\) of another monic polynomial \(f_0(t,X)\)
regular at \(a\) with separable reduction. Under the valuation
induced from \(L\), the residue field \(\ell_0\) of \(L_0\) is
the splitting field of \(f_0(a,X)\) inside \(\ell\).
Under the decomposition group isomorphisms of
Theorem~\ref{res:IG10}, restriction from the subgroup for \(L\)
to the subgroup for \(L_0\) is exactly the surjective map
\(\operatorname{Gal}(\ell/k)\to\operatorname{Gal}(\ell_0/k)\).
\end{corollary}

\begin{proof}
Every root of \(f_0(a+u,X)\) in \(L_0\) is integral for the
series valuation. If a root had negative order, its leading monic
term in the polynomial equation would have strictly lower order
than each of the other terms, which cannot sum to zero.
Reducing the monic factorization gives all the roots of \(f_0(a,X)\)
in \(\ell\), and they are distinct by separability. Let \(\ell_0\)
be the field they generate. The same roots lift in
\(\ell_0[[u]]\) by Lemma~\ref{res:IG09}, and uniqueness of the
lifts inside \(\ell[[u]]\) identifies these with the original
roots. Consequently \(L_0\) embeds into \(\ell_0((u))\), and its
residue field is exactly \(\ell_0\). Reduction commutes with
restriction to \(L_0\). The residue restriction is surjective
by finite Galois restriction, and the two vertical identifications
are those of Theorem~\ref{res:IG10}.
\end{proof}

\section{Regular realizations and arithmetic examples}

A finite Galois realization \(L/k(t)\) is called regular over
\(k\) when \(k\) is algebraically closed in \(L\): every element
of \(L\) algebraic over \(k\) lies in \(k\).
Here the finite extension of \(k(t)\) is separable already.
The criterion below works over imperfect fields as well.

\begin{proposition}[Regularity and the geometric group]
\label{res:IG12}
For a finite Galois extension \(L/k(t)\), work in a common algebraic
closure of \(k(t)\) and write \(\overline k\) for its algebraic
closure of constants. Restriction embeds
\(\operatorname{Gal}(L\overline k(t)/\overline k(t))\) into
\(\operatorname{Gal}(L/k(t))\). It is onto exactly when \(L/k(t)\)
is regular over \(k\).
\end{proposition}

\begin{proof}
Let \(M=L\cap\overline k(t)\). Restriction identifies the geometric
group with \(\operatorname{Gal}(L/M)\) by Galois base change
\cite[Result C10]{mg-aft}. We show \(M=C(t)\) for a finite
extension \(C/k\) inside \(L\cap\overline k\).

In characteristic \(p>0\), each member of \(\overline k(t)\)
has some \(p\)-power in \(k^{\mathrm{s}}(t)\): take an exponent
making all its finitely many algebraic coefficients separable over
\(k\). For an algebraic constant \(\beta\), write its monic minimal
polynomial as \(g(X^{p^r})\) with \(g'\ne0\) by collecting successive
\(p\)-th powers in the exponents. Irreducibility of \(g\) follows
from that of the original polynomial, so \(\beta^{p^r}\) is
separable. Any \(x\in M\) is separable over \(k(t)\), and hence
separable over \(k^{\mathrm{s}}(t)\), since its minimal polynomial
over the larger field divides a separable polynomial. It is also
purely inseparable over \(k^{\mathrm{s}}(t)\), so it belongs there.
In characteristic zero \(k^{\mathrm{s}}=\overline k\) already.
Since \(M/k(t)\) is finite, its finitely many generators have
their coefficients in some finite Galois extension \(k'/k\)
inside \(k^{\mathrm{s}}\), using Result F03 of
\cite{mg-aft}.

The extension \(k'(t)/k(t)\) has group \(\operatorname{Gal}(k'/k)\)
acting on coefficients. Indeed any constant basis over \(k\)
remains independent over \(k(t)\) by clearing denominators and
comparing powers of \(t\). The resulting span over \(k(t)\)
is a finite-dimensional domain containing \(k'\) and \(k(t)\).
Multiplication by a nonzero element is injective and hence
surjective, so this span is a field and equals \(k'(t)\).
Thus the coefficient actions already supply as many automorphisms
as the extension degree. Let \(H\) be the subgroup fixing \(M\),
and put \(C=(k')^H\). Then \(C(t)\subseteq M\); the same
basis argument gives \([k'(t):C(t)]=[k':C]=|H|\), which is
also \([k'(t):M]\) by finite Galois correspondence. Hence
\(M=C(t)\), with \(C\subseteq L\cap\overline k\).
If \(L\cap\overline k=k\), this forces \(M=k(t)\).
Conversely if \(M=k(t)\), an algebraic constant in \(L\) is
a constant rational function in \(k(t)\), thus lies in \(k\)
by comparison of coefficients after clearing the denominator.
\end{proof}

\begin{theorem}[A regular \(S_3\)-family over \(\mathbb Q\)]
\label{res:IG13}
The splitting field of \(X^3-X-t\) over \(\mathbb Q(t)\)
is regular over \(\mathbb Q\) with group \(S_3\).
Every odd integer \(a\) gives a good specialization with
Galois group \(S_3\) over \(\mathbb Q\).
\end{theorem}

\begin{proof}
There is no root in \(\overline{\mathbb Q}(t)\).
For a putative root \(P(t)/Q(t)\) in lowest terms, the equation
\(P^3-PQ^2=tQ^3\) forces \(Q\mid P^3\) in
\(\overline{\mathbb Q}[t]\). Thus \(Q\) is constant, and
a nonconstant polynomial \(h\) with \(h^3-h=t\) is impossible
by degree. A reducible cubic over a field would have a root.

For a monic cubic with roots \(r_1,r_2,r_3\) counted with multiplicity,
a square root of its discriminant is
\(d=\prod_{i<j}(r_i-r_j)\). Permuting the roots acts on \(d\)
by the sign. For \(f=X^3-X-t\) one has
\[
 d^2=-\prod_i f'(r_i)
     =-27f(b)f(-b)=4-27t^2,\qquad b^2=1/3 .
\]
The first equality follows by taking each ordered pair of
distinct indices in the derivative product; the second follows
from \(f'(X)=3(X-b)(X+b)\). The displayed polynomial is
nonzero, so the cubic is separable, and it has two distinct
simple zeros in \(\overline{\mathbb Q}\), so it is no square
in \(\overline{\mathbb Q}(t)\). For an irreducible separable
cubic in characteristic zero, its Galois group is a transitive
subgroup of \(S_3\). It has order divisible by three and it
lies in \(A_3\) exactly when \(d\) lies in the ground field.
Indeed this is the fixed-field criterion for \(d\), and
\(d^2\) is a square in the ground field exactly when \(d\) is
in it. Thus a nonsquare discriminant gives order six.
This proves that the geometric group is \(S_3\); the arithmetic
group also acts faithfully on the three roots and contains the
geometric group. Regularity follows by Proposition~\ref{res:IG12}.

If \(a\) is odd, the specialization \(X^3-X-a\) reduces to
\(X^3+X+1\) modulo \(2\), which has neither \(0\) nor \(1\)
as a root. A rational root of a monic integer cubic is an
integer, as its denominator in lowest terms divides the cube
of its numerator. There is no such root here. Its discriminant
\(4-27a^2\) is negative, so it is nonzero and nonsquare over
\(\mathbb Q\). The same transitivity and sign argument gives
group \(S_3\). Theorem~\ref{res:IG10} identifies this specialized
group with the full generic group.
\end{proof}

\begin{example}[A regular cyclic family]
\label{res:IG14}
Let \(n\geq2\) be invertible in \(k\), and suppose \(k\) contains a
primitive \(n\)-th root of unity \(\zeta\). The
splitting field of \(X^n-t\) over \(k(t)\) is a regular realization
of \(C_n\). The exponent of the factor \(t\) in a rational
function, with poles counted negatively, is additive under
multiplication. Hence \(t^i\) can be an \(n\)-th power in \(k(t)\)
only if \(n\mid i\), and the class of \(t\) has order \(n\).
Kummer duality (Theorem~\ref{res:GH07}) gives the cyclic group.
The same exponent argument over \(\overline k(t)\) gives the
full geometric group, so regularity follows from
Proposition~\ref{res:IG12}. For instance \(n=2\) over
\(\mathbb Q\) gives full group at each positive prime \(t=p\)
by Theorem~\ref{res:IG10}: the specialization has distinct roots
and \(p\) is nonsquare by prime factorization.
\end{example}

\begin{lemma}[Fresh primes for prescribed cyclic factors]
\label{res:IG15}
Given \(n\geq2\) and a finite set of primes, there is a prime
\(p\) outside the set with \(n\mid p-1\).
\end{lemma}

\begin{proof}
Use the monic integral cyclotomic polynomials and their product
formula from \cite[Result K07.15.cyclotomic\_polynomials]{mg-d07}.
In that formula \(\Phi_1(0)=-1\); evaluation at zero and induction
on \(n\) give \(\Phi_n(0)=1\) for \(n>1\).
Choose an integer \(b\geq3\) divisible by \(n\) and all the
excluded primes. Since each primitive complex \(n\)-th root
\(\xi\) has modulus \(1\), each factor \(|b-\xi|\) is at least
\(b-1\). Thus \(|\Phi_n(b)|>1\). Choose a prime divisor \(p\)
of \(\Phi_n(b)\); it does not divide \(b\), since
\(\Phi_n(b)\equiv 1\) modulo every prime dividing \(b\).
In particular \(p\nmid n\). Modulo \(p\) the polynomial
\(X^n-1\) has no repeated root, since its derivative is
\(nX^{n-1}\). The factors \(\Phi_d\) for \(d\mid n\) therefore
have disjoint root sets. The class of \(b\) has \(n\)-th power
one; if its order were \(d<n\) with \(d\mid n\), it would
also be a root of a factor of \(X^d-1\), a contradiction.
Hence it has order \(n\) in \(\mathbb F_p^\times\), and
\(n\mid p-1\).
\end{proof}

\begin{theorem}[Abelian realizations disjoint from a prescribed field]
\label{res:IG16}
Every finite abelian group \(A\) occurs as a Galois group over
\(\mathbb Q\). More precisely, for any finite Galois extension
\(E/\mathbb Q\) there is a finite Galois extension \(F/\mathbb Q\)
with group \(A\) and \(F\cap E=\mathbb Q\). Consequently the
split embedding problem
\[
   A\times\operatorname{Gal}(E/\mathbb Q)
             \longrightarrow \operatorname{Gal}(E/\mathbb Q)
\]
over the restriction to \(E\) has a proper solution through \(EF\).
\end{theorem}

\begin{proof}
Write \(A=\prod_{i=1}^r C_{n_i}\) with \(n_i\geq2\),
using \cite[Result K07.07.invariant\_factor\_existence]{mg-d07}.
The case \(r=0\) takes \(F=\mathbb Q\). Choose any finite
number \(s\) of disjoint prime lists \(p_{j,i}\), with
\(n_i\mid p_{j,i}-1\), by Lemma~\ref{res:IG15}.
Let \(m\) be the product of all selected primes. The cyclotomic
group of \(\mathbb Q(\zeta_m)\) is
\((\mathbb Z/m\mathbb Z)^\times\) by
\cite[Result K07.15.cyclotomic\_galois]{mg-d07}.
By the Chinese remainder isomorphism it is the product of
the groups \(\mathbb F_{p_{j,i}}^\times\), each cyclic of
order \(p_{j,i}-1\) by
\cite[Result K07.13.multiplicative\_cyclic]{mg-d07}.
For list \(j\), project to its factors and take their respective
quotients of order \(n_i\), giving a surjection onto a copy
of \(A\). The fixed field \(F_j\) of its kernel is Galois
with group \(A\). For two distinct lists the joint map onto
\(A\times A\) is surjective; thus
\([F_jF_h:\mathbb Q]=|A|^2\) for \(j\ne h\), by the
fixed-field correspondence. Galois base change
\cite[Result C10]{mg-aft} now gives \(F_j\cap F_h=\mathbb Q\).

The finite Galois extension \(E/\mathbb Q\) has only finitely
many intermediate fields, since its Galois group has finitely
many subgroups. Take \(s\) greater than this number. If all
intersections \(F_j\cap E\) were nontrivial over \(\mathbb Q\),
they would be pairwise different because the \(F_j\) have
trivial pairwise intersection. This is impossible, so some
\(F=F_j\) has \(F\cap E=\mathbb Q\). Restriction from
\(\operatorname{Gal}(EF/\mathbb Q)\) to the product for \(F\)
and \(E\) is injective since the two fields generate \(EF\);
it is surjective by the degree formula from Result C10.
The resulting isomorphism has second projection equal to
restriction to \(E\). Proposition~\ref{res:IG02} gives the
proper solution.
\end{proof}

